% Fragment: \input z poziomu tgp_yukawa_exact_reduction.tex
% Wybrany wariant P4 (plan): rozwinięcie multipolowe = szereg Legendre'a / Gegenbauera
% dla jądra Yukawy, redukcja I_Y do całek radialnych + całek kątowych (Gaunt).

%==========================================================================
\section{Wybrana \'{s}cie\.{z}ka P4: multipolowo\'s\'c (Legendre / Gegenbauer)}
\label{sec:p4_multipole_gegenbauer}
%==========================================================================

Zamiast poszukiwa\'c sko\'nczonej kombinacji funkcji Bessela (por.
twierdzenie~\ref{thm:nonelementary} poni\.{z}ej), naturalnym \textbf{kontynuacj\k{a}
wyprowadzenia} $I_Y$ jest rozwini\k{e}cie kuliste j\k{a}dra
$G_m(R)=e^{-mR}/R$ w szereg wielomian\'ow Legendre'a --- tj.\ rozwini\k{e}cie
typu Gegenbauera $C_\ell^{(1/2)}=P_\ell$ dla zielonej funkcji operatora
Helmholtza w $\mathbb{R}^3$.

\subsection{Twierdzenie dodawania (trzy wymiary)}

Niech $i_\ell,\,k_\ell$ oznaczaj\k{a} \emph{sferyczne} zmodyfikowane funkcje Bessela
(DLMF~10.47; w Pythonie \texttt{scipy.special.spherical\_in},
\texttt{spherical\_kn}). Dla $r=|\mathbf{r}|$, $r'=|\mathbf{r}'|$,
$r_<\!=\min(r,r')$, $r_>\!=\max(r,r')$ oraz
$\cos\gamma=(\mathbf{r}\cdot\mathbf{r}')/(rr')$ zachodzi klasyczna reprezentacja
\begin{equation}
\label{eq:addition_yukawa}
G_m\bigl(|\mathbf{r}-\mathbf{r}'|\bigr)
=
\sum_{\ell=0}^{\infty} (2\ell+1)\,
\mathsf{w}_\ell(m,r,r')\,
P_\ell(\cos\gamma),
\end{equation}
gdzie wsp\'o\l{}czynnik radialny ma posta\'c iloczynu fal cz\k{a}stkowych Yukawy:
\begin{equation}
\label{eq:w_l_yukawa}
\mathsf{w}_\ell(m,r,r')
\;\equiv\;
\frac{2m}{\pi}\,
i_\ell(m r_<)\, k_\ell(m r_>).
\end{equation}

\begin{remark}[Konwencja i sp\'ojno\'s\'c]
Wsp\'o\l{}czynnik \eqref{eq:w_l_yukawa} jest standardowy przy powy\.{z}szej
normalizacji $i_\ell,k_\ell$ jak w DLMF; w praktyce runtime nale\.{z}y
\textbf{zweryfikowa\'c numerycznie} zgodno\'s\'c \eqref{eq:addition_yukawa} z
reprezentacj\k{a} Schwingera \eqref{eq:IY_schwinger_2D} na kilku parach
$(r,r',\gamma)$ zanim u\.{z}yje si\k{e} szeregu do przybli\.{z}ania $I_Y$.
Granica $m\to0$ powinna odtwarza\'c rozwini\k{e}cie Laplace'a
$\sum_\ell r_<^\ell r_>^{-\ell-1}P_\ell$ (test prefaktora).
\end{remark}

\subsection{Jeden \'{s}rodek rozwini\k{e}cia: $\mathbf{x}_3=\mathbf{0}$}

Obieramy uk\l{}ad z $\mathbf{x}_3=\mathbf{0}$, $|\mathbf{x}_1|=d_{13}$,
$|\mathbf{x}_2|=d_{23}$, a $\omega$ --- k\k{a}tem mi\k{e}dzy promieniami
$\mathbf{x}_1,\mathbf{x}_2$ wzgl\k{e}dem $\mathbf{x}_3$:
\[
\cos\omega=\frac{d_{13}^2+d_{23}^2-d_{12}^2}{2d_{13}d_{23}}.
\]
Wtedy
\begin{align}
G_m(|\mathbf{x}-\mathbf{x}_1|)
&=
\sum_{\ell=0}^{\infty}(2\ell+1)\,
\mathsf{w}_\ell(m,r,d_{13})\,
P_\ell\bigl(\hat{\mathbf{x}}\cdot\widehat{\mathbf{x}_1}\bigr),
\\
G_m(|\mathbf{x}-\mathbf{x}_2|)
&=
\sum_{\ell'=0}^{\infty}(2\ell'+1)\,
\mathsf{w}_{\ell'}(m,r,d_{23})\,
P_{\ell'}\bigl(\hat{\mathbf{x}}\cdot\widehat{\mathbf{x}_2}\bigr),
\\
G_m(|\mathbf{x}-\mathbf{x}_3|) &= G_m(r)=\frac{e^{-mr}}{r},
\qquad r=|\mathbf{x}|.
\end{align}

Ca\l{}ka po k\k{a}cie w
$I_Y=\int d^3x\,G_m(r)\,G_m(|\mathbf{x}-\mathbf{x}_1|)\,G_m(|\mathbf{x}-\mathbf{x}_2|)$
sprowadza si\k{e} do ca\l{}ek postaci
\begin{equation}
\label{eq:angular_A}
\mathcal{A}_{\ell\ell'}(\omega)
\;\equiv\;
\int_{S^2}
P_\ell\bigl(\hat{\mathbf{x}}\cdot\hat{\mathbf{n}}_1\bigr)\,
P_{\ell'}\bigl(\hat{\mathbf{x}}\cdot\hat{\mathbf{n}}_2\bigr)\,d\Omega,
\qquad
\hat{\mathbf{n}}_1\cdot\hat{\mathbf{n}}_2=\cos\omega,
\end{equation}
kt\'ore s\k{a} \textbf{wielomianami} zmiennej $\cos\omega$ (kombinacje symboli
$3j$ / wsp\'o\l{}czynnik\'ow Gaunta dla zonalnych harmonic sferycznych).
Dla $\ell\neq\ell'$ ca\l{}ka \eqref{eq:angular_A} nie znika w og\'olno\'sci ---
pojawia si\k{e} pe\l{}na podw\'ojna suma po $(\ell,\ell')$.

\subsection{Szkielet radialny}

Po rozdzieleniu $d^3x=r^2\,dr\,d\Omega$ otrzymujemy
\begin{equation}
\label{eq:IY_multipole_double_sum}
\boxed{
I_Y
=
\sum_{\ell=0}^{\infty}\sum_{\ell'=0}^{\infty}
(2\ell+1)(2\ell'+1)\,
\int_{0}^{\infty} dr\; r^2\,\frac{e^{-mr}}{r}\,
\mathsf{w}_\ell(m,r,d_{13})\,
\mathsf{w}_{\ell'}(m,r,d_{23})\,
\mathcal{A}_{\ell\ell'}(\omega).
}
\end{equation}
Wielko\'s\'c $I_Y$ zale\.{z}y tylko od $(d_{12},d_{13},d_{23},m)$, wi\k{e}c
$\omega$ wyra\.{z}a si\k{e} przez prawo cosinus\'ow jak wy\.{z}ej; symetria
permutacyjna tripletu mo\.{z}e by\'c odzyskana przez u\'srednienie lub wyb\'or
innego centrum rozwini\k{e}cia.

\subsection{Co daje ten wariant w praktyce (runtime)}

\begin{itemize}
\item \textbf{Semianalityka.} Uci\k{e}cie $\ell,\ell'\le L_{\max}$ zamienia 2D
simpleks z \texttt{three\_body\_force\_exact.py} na ci\k{a}g 1D ca\l{}ek
radialnych (kwadratura Gaussa na $r$) z \textbf{sta\l{}\k{a}} liczb\k{a}
wyraz\'ow podw\'ojnej sumy $O(L_{\max}^2)$ --- dogodne do GPU / tabulacji
$(d_{ij},m)$ po wst\k{e}pnym wyliczeniu $\mathcal{A}_{\ell\ell'}$ jako
wielomian\'ow w $\cos\omega$.
\item \textbf{Zwi\k{a}zek z nieelementarno\'sci\k{a}.} Szereg
\eqref{eq:IY_multipole_double_sum} jest co do zasady \textbf{niesko\'nczony};
nie zaprzecza twierdzeniu~\ref{thm:nonelementary} o braku sko\'nczonej
kombinacji $K_\nu(a_k t)$ dla $f_\triangle(t)$ --- pokazuje natomiast, \.{z}e
$I_Y$ jest naturalnie reprezentowalne jako granica sum cz\k{e}\'{s}ciowych
fal multipolowych.
\item \textbf{Nast\k{e}pny krok techniczny (do kodu).} Wypisa\'c jawnie
$\mathcal{A}_{\ell\ell'}(\omega)$ dla ma\l{}ych $L_{\max}$ (np.\ $2$ lub $3$)
oraz por\'owna\'c \eqref{eq:IY_multipole_double_sum} z
\texttt{yukawa\_overlap\_exact} na siatce $(d_{12},d_{13},d_{23})$ --- to
zamknie p\k{e}tl\k{e} weryfikacyjn\k{a} dla prefaktora \eqref{eq:w_l_yukawa}.
\end{itemize}
